\chapter{Addendum: The Debugging Agent}
\label{sec:addendum}

This chapter describes work added on 24 November 2026, after this report was
submitted. It is not part of the submitted evaluation: the results of
Chapter~\ref{sec:results}, the test counts of Table~\ref{tab:verify} and the
conclusions of Chapter~\ref{sec:conclusion} are those of the submitted version and do
not depend on anything here. The chapter is included so that the feature can be
examined with the same standard of evidence as the rest of the report. The code is on
the branch \code{agent-debug}, commit \code{3d05b94}.

\section{Motivation}
\label{sec:agentmotiv}
The original SiliconScribe is a fixed pipeline (Chapter~\ref{sec:arch}). The
orchestrator decides what runs and in which order, and the model fills each step:
the specification, the RTL, the testbench and, on a failure, the diagnosis. This has
two consequences. First, the model never chooses what to inspect. The fixer receives
the RTL, the testbench, the specification and the simulator log
(Section~\ref{sec:loop}) and must diagnose from those alone; it cannot ask for the
value of a signal at a given time or for the clock cycles before a failing check.
Second, the user cannot ask a follow-up question about a run. A failing run shows the
failing checks and the waveform, but connecting a wrong value to its cause is left to
the user.

The debugging agent, shown in the interface as ``Ask about this run'', adds the
missing interaction. The model chooses which tools to call, but every tool result
comes from the simulator or from the source code, and every fix it proposes is
simulated again by the server. The rule of the rest of the system, that only the
simulator decides PASS or FAIL, still holds.

\section{Design}
\label{sec:agentdesign}
One question is handled as follows (Figure~\ref{fig:agentflow}):

\begin{enumerate}
\item \textbf{Question.} The user types a question about the code currently in the
editor, or clicks one of the examples (``Why did the first check fail?'', ``Is the
design or the testbench wrong?'', ``Walk me through the first 5 clock cycles.'').
The last four questions and answers of the conversation are sent as context.
\item \textbf{Evidence run.} The server simulates the editor's code again in the
same sandbox as Re-run (Section~\ref{sec:sandbox}): Icarus Verilog, allow-listed
system tasks, CPU, memory and file limits, and a 30~s limit. The answer is therefore
about exactly the code on screen. If the code was edited since the last run, the
Results and Waveform tabs switch to this new run.
\item \textbf{Overview.} The model receives a short overview: the verdict, the first
failing checks, the clocks and their period, the list of signals in the waveform and
the numbered RTL. The testbench is not included; the model reads it with
\code{read_code} when it needs it.
\item \textbf{Tool loop.} Each model reply is one JSON object, either a tool call or
the final answer. The tool result is appended to the transcript and the model is
called again. At most 10 model calls are made per question.
\item \textbf{Answer.} The final answer contains markdown text, a suspect (design,
testbench, both, unclear or none), a confidence (high, medium or low) and, optionally,
a fix as a complete corrected file. A fix is simulated by the server before it is
shown.
\end{enumerate}

The server streams the progress to the browser as Server-Sent Events: a status
message, the evidence run with its waveform, one event per tool call and then exactly
one answer (or an error).

\begin{figure}[H]
\centering
\begin{tikzpicture}[
    font=\footnotesize,
    start/.style={draw=ink,rounded corners=8pt,fill=softgrey,text width=2.6cm,align=center,inner sep=4pt},
    proc/.style={draw=ink,rounded corners=2pt,fill=white,text width=2.6cm,align=center,inner sep=4pt},
    dec/.style={draw=ink,diamond,aspect=2.2,fill=softgrey,text width=1.9cm,align=center,inner sep=1pt},
    arr/.style={-{Latex[length=2.4mm]},line width=0.7pt},
    lbl/.style={font=\scriptsize,color=inkdark}
  ]
  \node[start] (q) {Question about the code in the editor};
  \node[proc,right=7mm of q] (run) {Evidence run\\\scriptsize sandbox, 30~s limit};
  \node[proc,right=7mm of run] (ov) {Overview\\\scriptsize verdict, failing checks, clocks, signals, RTL};
  \node[dec,below=8mm of ov] (m) {Model reply: tool or answer?};
  \node[proc,left=14mm of m] (t) {Tool\\\scriptsize waveform, source or \code{test_fix}};
  \node[proc,below=8mm of m] (fix) {Server re-simulates the fix};
  \node[start,left=14mm of fix] (ui) {Answer, suspect, confidence, fix verdict};
  \draw[arr] (q) -- (run);
  \draw[arr] (run) -- (ov);
  \draw[arr] (ov) -- (m);
  \draw[arr] (m.west) -- node[lbl,above]{tool} (t.east);
  \draw[arr] (t.north) -- ++(0,0.5) -| node[lbl,above,pos=0.25]{result} (m.north west);
  \draw[arr] (m) -- node[lbl,right]{answer} (fix);
  \draw[arr] (fix) -- (ui);
\end{tikzpicture}
\caption[Flow of one question to the debugging agent]{Flow of one question to the
debugging agent. The model chooses the tools; every tool result comes from the
simulation or the source, and a proposed fix is simulated again before it is shown.
At most 10 model calls and 3 fix tests are made per question.}
\label{fig:agentflow}
\end{figure}

\subsection{Tools}
Table~\ref{tab:agenttools} lists the eight tools. Times are given in nanoseconds.
Five of them read the simulation (\code{failing_checks}, \code{values_at},
\code{changes}, \code{cycle_table}, \code{find_value}), two read the source
(\code{drivers}, \code{read_code}) and one runs a new simulation (\code{test_fix}).
The cycle table samples each signal just before the clock edge, which is the value
the flip-flops see, and just after it; this is the distinction that matters when
reading sequential logic from a waveform.

\begin{table}[H]
\centering
\caption{The eight tools available to the debugging agent.}
\label{tab:agenttools}
\footnotesize
\begin{tabularx}{\textwidth}{@{}p{2.4cm}X@{}}
\toprule
\textbf{Tool} & \textbf{What it returns} \\
\midrule
\code{failing_checks} & The testbench's failed comparisons: time, name, expected value and actual value. \\
\code{values_at} & The value of each requested signal at a given time, and just before it. \\
\code{changes} & Every value change of one signal in a time window (at most 40). \\
\code{cycle_table} & One row per clock edge, with each signal sampled just before the edge (what the flip-flops see) and just after it. \\
\code{find_value} & The times at which a signal held a given value (decimal, hexadecimal, binary, \code{x} or \code{z}). \\
\code{drivers} & The source lines that assign a signal, with the enclosing \code{always} or \code{assign} block, and the instance ports connected to it. This is a text search that ignores comments and comparisons, not a full elaboration. \\
\code{read_code} & Numbered source lines of the RTL or the testbench (at most 120 lines). \\
\code{test_fix} & Simulates a candidate fix (a complete file) with the other file unchanged and returns the verdict. At most 3 per question. \\
\bottomrule
\end{tabularx}
\end{table}

\subsection{Protocol}
Each model reply must be exactly one JSON object of one of the two forms below
(shortened):

\begin{lstlisting}[language=,basicstyle=\ttfamily\scriptsize]
{"thought": "<what is checked and why>", "tool": "cycle_table",
 "args": {"signals": ["enable", "count"], "from_ns": 0, "to_ns": 100}}

{"answer": "<markdown>", "suspect": "design", "confidence": "high",
 "fix": {"file": "rtl", "code": "<the complete corrected file>"}}
\end{lstlisting}

The protocol is plain JSON text, not a provider's native function-calling interface.
Native function calling differs between providers and is not available for every
model, whereas any chat model can be asked to reply with a JSON object. The agent therefore
works with every model source the application supports: the OpenCode app,
Cloudflare Workers AI, NVIDIA NIM and any OpenAI-compatible endpoint. The fixer of
Section~\ref{sec:loop} moved away from JSON after a live run produced malformed JSON
with a complete Verilog file inside it. The agent's final answer can also carry a
complete file inside JSON, so the same risk exists; most of its replies, however, are
short tool calls, and a reply that is not one JSON object is handled explicitly
(Section~\ref{sec:agenttrust}).

Each step sends the whole transcript again. To limit its growth, one tool result is
cut to 3000 characters, and results older than the last five steps are shortened to
400 characters.

\section{Trust Rules}
\label{sec:agenttrust}
The agent follows the same policy as the rest of SiliconScribe: the simulator
decides, and anything that the simulator did not check is labelled as such.

\begin{enumerate}
\item \textbf{Grounding rule.} An answer must rest on at least one simulation query
(\code{values_at}, \code{changes}, \code{cycle_table}, \code{find_value} or
\code{test_fix}). If the model answers without one, it is sent back once to check;
a second answer is accepted. The rule was added after live case A4
(Section~\ref{sec:agentlive}).
\item \textbf{Every fix is re-simulated by the server.} The fix in the final answer
is simulated again before it is shown, and the interface shows that verdict (for
example ``PASS, 57 of 57 checks''), never the model's own claim. A ``fix'' identical
to the current file is dropped, and a fix that cannot be simulated is shown as an
error. The fix is not applied automatically: ``Apply to the editor'' places it in the
editor, and Re-run then shows it in the main results.
\item \textbf{A testbench fix is flagged.} A fix to the testbench is shown with the
warning ``This changes the checker itself. A testbench fix can make a wrong design
pass.''
\item \textbf{Labelled rule-based fallback.} When no AI model is connected (offline),
when the model fails (quota, timeout or error) or when the 10-step limit is reached,
the panel shows rule-based evidence instead: the first failing check, the port values
at that time, the clock cycles before it and the lines that assign the outputs. It is
labelled ``Rule-based evidence (no AI)'' and ``not an AI diagnosis''. As in
Section~\ref{sec:fallback}, there is never a silent swap.
\item \textbf{Unchecked text is marked.} If the model twice ignores the JSON
protocol, its text is shown with the note that this text is unchecked.
\item \textbf{No HTML from the model.} Model text is rendered as plain text, with
only code, bold and times formatted; it is never inserted as HTML.
\item \textbf{The sandbox applies throughout.} The policy of
Section~\ref{sec:sandbox} applies to the evidence run and to every fix test; a
\code{$system} call is refused.
\end{enumerate}

What is \emph{not} verified is the explanation text itself. Only the fix is checked
by simulation. Live case A6 below shows why this distinction matters.

\section{User Interface}
\label{sec:agentui}
The panel ``Ask about this run'' sits in the left sidebar, under the run timeline
(Figure~\ref{fig:askpanel}). The evidence steps appear as they happen, followed by
the answer, the suspect and confidence, and a fix card with the server's verdict.
Times written in an answer (``@25ns'' or ``25ns'') are buttons: clicking one opens
the Waveform tab with cursor A at that time (Figure~\ref{fig:askwave}).
Figure~\ref{fig:askanswer} shows the whole window. The three screenshots were taken
from the real live run A3 (light theme); nothing in them is mocked.

\begin{figure}[htbp]
\centering
\includegraphics[height=11cm]{s22_ask_panel.png}
\caption[The ``Ask about this run'' panel]{The whole ``Ask about this run'' panel in
live run A3: the question, the evidence steps, the answer and the fix card.}
\label{fig:askpanel}
\end{figure}

\begin{figure}[htbp]
\centering
\includegraphics[width=0.85\textwidth]{s23_ask_waveform_jump.png}
\caption[Jumping from a cited time to the waveform]{After clicking ``@30ns'' in the
answer: the Waveform tab with cursor A at 30\,ns, where \texttt{enable} = 0 and
\texttt{count} = 0x2 while the hold check expected 1.}
\label{fig:askwave}
\end{figure}

\begin{figure}[htbp]
\centering
\includegraphics[width=0.85\textwidth]{s21_ask_answer.png}
\caption[Answer and main results of the edited code]{The whole window in live run A3:
the answer in the sidebar (suspect Design, cited times, fix PASS 57/57) and the
Results tab showing the edited code's FAIL 12/57 (45 failed).}
\label{fig:askanswer}
\end{figure}

\section{Live Results}
\label{sec:agentlive}
Table~\ref{tab:agentlive} lists the live runs made on 24 November 2026. Each case was
run once, on small designs, and the model was NVIDIA NIM
\code{nvidia/nemotron-3-ultra-550b-a55b} unless stated.

\begin{table}[H]
\centering
\caption[Live runs of the debugging agent]{Live runs of the debugging agent,
24 November 2026 (one run each). Time is the time to the answer.}
\label{tab:agentlive}
\footnotesize
\begin{tabularx}{\textwidth}{@{}l >{\raggedright\arraybackslash}p{4.3cm} >{\raggedright\arraybackslash}X l@{}}
\toprule
\textbf{Case} & \textbf{Set-up} & \textbf{Outcome} & \textbf{Time} \\
\midrule
A1 & Built-in buggy counter (adds 2; the RTL has a ``BUG'' comment) & Correct, suspect design; fix re-simulated 57/57 & 31~s \\
A2 & Hidden bug \code{else if (enable || count[0])}, no comment & Correct; used a cycle table (\code{count} moves while \code{enable}=0 at 25, 55 and 85\,ns); fix 57/57 & 34~s \\
A3 & Same hidden bug, in the real browser interface (twice) & Correct both times, fix 57/57; one answer cited 9 clickable times & 107~s, 74~s \\
A4 & Same hidden bug, model \texttt{opencode/\allowbreak deepseek-v4.1-flash-free} & Correct, fix 57/57, but zero tool calls (reasoned from the overview) & 49~s \\
A5 & Correct counter, wrong testbench (expects 3 instead of 2 in \code{counting_continues}) & Suspect testbench: the check runs at the negative edge at 560\,ns, before the third positive edge; testbench fix 57/57 & 219~s, 5 steps \\
A6 & Saved SR-latch project that timed out (a real AI design from 20 November), real browser & Fix (gate delays) re-simulated 558/558, but the explanation was only partly right (see text) & 316~s \\
A7 & Older saved counter project whose 20-check testbench never drops \code{enable}, with the hidden bug & Reported that the run passed and that no hold checks failed, instead of inventing a failure & about 170~s \\
\bottomrule
\end{tabularx}
\end{table}

\paragraph{The hidden bug (A1 to A4).} A1 is the easy case: the RTL carries a ``BUG''
comment. In A2 the bug, \code{else if (enable || count[0])}, has no comment, and the
agent found it from the waveform: its cycle table showed \code{count} changing while
\code{enable} was 0 at 25, 55 and 85\,ns. A3 repeated the case twice in the real
browser interface, with the same result. In A4 a different model,
\code{opencode/deepseek-v4.1-flash-free}, gave the correct answer and a passing fix
without calling a single tool: it reasoned from the overview alone. The answer was
right, but nothing in it came from the run. This case led to the grounding rule of
Section~\ref{sec:agenttrust}.

\paragraph{The checker was wrong (A5).} In A5 the counter was correct and the
testbench was wrong: the check \code{counting_continues} expected 3 where 2 is
correct. The agent named the testbench as the suspect, with evidence: the check runs
at the negative edge at 560\,ns, before the third positive edge. Its testbench fix
was re-simulated at 57 of 57 checks. The agent can therefore say that the design is
fine and the checker is wrong, and show why. Because this fix changes the checker,
the interface shows it with the testbench warning, and the decision to apply it is
left to the user.

\paragraph{Fix passed, explanation partly wrong (A6).} A6 used a saved SR-latch
project, a real AI-generated design from 20 November whose run had timed out. The
agent's fix, which added gate delays, was re-simulated by the server and passed 558
of 558 checks. Its explanation, however, was only partly right: it blamed the hold
state, whereas the real cause, found by hand on 20 November, is that the testbench
releases both inputs from the forbidden 0,0 state at the same time. The fix is
verified; the explanation is model text and is not. This is the clearest example of
the limit stated at the end of Section~\ref{sec:agenttrust}. Because the design hangs,
the evidence run also used its full 30~s before the model started.

\paragraph{A false premise (A7).} In A7 the question assumed a failure that did not
happen: the older counter project still had the hidden bug, but its 20-check
testbench never drops \code{enable}. The agent reported that the run passed and
that no hold checks failed, instead of inventing a failure. This weak testbench
does not expose the bug.

In summary, in every case with a fix (A1 to A6) the fix was shown only with the
server's own re-simulation verdict. Answers took about 30~s to 5 minutes on free
models.

\section{Tests}
\label{sec:agenttests}
All tests pass on commit \code{3d05b94}:

\begin{itemize}
\item \textbf{Backend:} 632 pytest tests (591 before the addendum). The 41 new tests
in \code{tests/test_agent.py} cover value parsing, sampling before and after a clock
edge, the tools on a real failing run, the code search ignoring comparisons, the loop
with a scripted model, fix verification, a wrong fix shown as failing, the grounding
rule, the fallback when the model fails, the step limit, the API in offline mode, the
sandbox and request validation.
\item \textbf{Frontend:} 277 vitest tests (267 before). The 10 new tests in
\code{DebugChat.test.tsx} cover the streamed events, clickable times in VCD ticks,
the fix card's verdict and ``Apply'', the testbench-fix warning, the rule-based
label, errors and history, model text never rendered as HTML, and the busy state.
\item \textbf{Sabotage checks:} making the fix check always report PASS, and making
the ``just before'' sample the same as the ``after'' sample, made 4 backend tests
fail; a wrong conversion from nanoseconds to VCD ticks made a frontend test fail. All
were restored afterwards.
\end{itemize}

\section{Limitations}
\label{sec:agentlimits}
\begin{itemize}
\item \textbf{The explanation is not verified.} The explanation is model text; only
the fix is checked by simulation (see A6).
\item \textbf{Slow.} A question took 30~s to 5 minutes on free models, because each
step sends the growing context again.
\item \textbf{Shallow code analysis.} \code{drivers} is a text search, not an
elaborated netlist. Only Icarus is used, so there are no SVA assertions and no UVM.
\item \textbf{Not a benchmark.} Each case was run once, on small designs and mainly
with one model.
\item \textbf{Hanging designs are slow to diagnose.} The evidence run is limited to
30~s, so a design that hangs (A6) costs the full 30~s before the model starts.
\end{itemize}
